\section{Frameworks and infinitesimal rigidity}
\label{sec:frameworks}

A $d$-dimensional \emph{framework} on a vertex set $V$ is a placement
$p : V \to \R^{d}$. Given a graph $G : \mathrm{SimpleGraph}\,V$, the
\emph{rigidity map} at $p$ is the $\R$-linear map sending an
infinitesimal motion $p' \in \R^{d \cdot |V|}$ to the family of
edge-length derivatives
$e \mapsto \langle p_u - p_v,\, p'_u - p'_v \rangle$ indexed by the
edges $e = \{u, v\} \in E(G)$. Geometrically this is the differential
at $p$ of the squared-edge-length map (up to a factor of $2$); we
define it directly via the explicit entry formula to avoid the
differentiability machinery that the rank/kernel arguments downstream
do not need.

A framework is \emph{infinitesimally rigid} when the kernel of the
rigidity map has dimension at most $d(d+1)/2$ --- the dimension of
trivial Euclidean motions (translations + infinitesimal rotations).
The graph $G$ is \emph{generically rigid} in dimension $d$ when
\emph{some} placement is infinitesimally rigid. The headline result of
this section is the edge-count bound: every generically rigid graph in
dimension $d$ on $n$ vertices satisfies $d \cdot n \le |E(G)| + d(d+1)/2$.
Specializing at $d = 2$ recovers the necessary half of Laman's
theorem.

All definitions are stated for general dimension $d : \N$; the
Laman-relevant specialization to $d = 2$ happens at the Phase~5
boundary.

\subsection{Definitions}

\begin{definition}[Framework]
  \label{def:framework}
  \lean{SimpleGraph.Framework}
  \leanok
  A $d$-dimensional \emph{framework} on $V$ is a function
  $p : V \to \mathrm{EuclideanSpace}\;\R\,(\mathrm{Fin}\,d)$. In Lean,
  $\mathrm{Framework}\,V\,d$ is an \texttt{abbrev} for this Pi type, so
  the Pi $\R$-module and finite-dimensional instances propagate
  transparently.
\end{definition}

\begin{lemma}[Dimension of the framework space]
  \label{lem:framework-finrank}
  \lean{SimpleGraph.Framework.finrank}
  \leanok
  \uses{def:framework}
  If $V$ is finite, then
  $\dim_{\R}\,(\mathrm{Framework}\,V\,d) = |V| \cdot d$.
\end{lemma}
\begin{proof}
  \leanok
  Apply $\mathrm{Module.finrank\_pi\_fintype}$: the framework space is
  the product of $|V|$ copies of $\R^{d}$, each of dimension $d$.
\end{proof}

\begin{definition}[Rigidity map]
  \label{def:rigidityMap}
  \lean{SimpleGraph.RigidityMap}
  \leanok
  \uses{def:framework}
  For a graph $G$ on $V$ and a framework $p : V \to \R^{d}$, the
  \emph{rigidity map} at $p$ is the $\R$-linear map
  \[
    \mathrm{RigidityMap}\,G\,p :
    \mathrm{Framework}\,V\,d \;\to_{\R}\; \bigl(E(G) \to \R\bigr)
  \]
  with
  \[
    (\mathrm{RigidityMap}\,G\,p)(p')(\{u, v\})
    \;=\;
    \langle p_u - p_v,\; p'_u - p'_v \rangle.
  \]
  The formula is symmetric in $u, v$ (each factor flips sign under
  swap), so it descends from an unordered pair $\{u, v\}$ to a single
  scalar via $\mathrm{Sym2.lift}$. Linearity in $p'$ is built
  compositionally from $\mathrm{innerSL}$ and the projection-and-
  subtraction maps, then bundled across edges by $\mathrm{LinearMap.pi}$.
\end{definition}

\begin{lemma}[Rigidity-map entry formula]
  \label{lem:rigidityMap-apply}
  \lean{SimpleGraph.rigidityMap_apply}
  \leanok
  \uses{def:rigidityMap}
  For $u, v \in V$ with $\{u, v\} \in E(G)$ and $p, p' : V \to \R^{d}$,
  \[
    (\mathrm{RigidityMap}\,G\,p)(p')(\{u, v\})
    \;=\;
    \langle p_u - p_v,\; p'_u - p'_v \rangle.
  \]
\end{lemma}
\begin{proof}
  \leanok
  Definitional, by $\mathrm{Sym2.lift}$ reduction on the
  representative pair.
\end{proof}

\begin{definition}[Infinitesimally rigid framework]
  \label{def:isInfinitesimallyRigid}
  \lean{SimpleGraph.IsInfinitesimallyRigid}
  \leanok
  \uses{def:rigidityMap}
  For finite $V$, a framework $p$ is \emph{infinitesimally rigid} for
  $G$ if
  \[
    \dim_{\R}\,\bigl(\ker (\mathrm{RigidityMap}\,G\,p)\bigr)
    \;\le\;
    \frac{d(d+1)}{2}.
  \]
  The bound is the dimension of the space of trivial Euclidean
  motions (translations + infinitesimal rotations), so this is the
  always-correct upper bound formulation: equality holds exactly when
  the kernel is exhausted by trivial motions.
\end{definition}

\begin{definition}[Generically rigid graph]
  \label{def:isGenericallyRigid}
  \lean{SimpleGraph.IsGenericallyRigid}
  \leanok
  \uses{def:isInfinitesimallyRigid}
  A graph $G$ on finite $V$ is \emph{generically rigid} in dimension
  $d$ if some framework $p : V \to \R^{d}$ is infinitesimally rigid
  for $G$.
\end{definition}

\subsection{Rigidity-map API}

\begin{lemma}[Kernel monotonicity]
  \label{lem:rigidityMap-ker-mono}
  \lean{SimpleGraph.rigidityMap_ker_mono}
  \leanok
  \uses{def:rigidityMap}
  If $G \le G'$ then for every framework $p$,
  \[
    \ker (\mathrm{RigidityMap}\,G'\,p)
    \;\le\;
    \ker (\mathrm{RigidityMap}\,G\,p).
  \]
\end{lemma}
\begin{proof}
  \leanok
  Adding edges only adds entries to the rigidity map; if $p'$ kills
  every edge of $G'$ it certainly kills every edge of $G \subseteq G'$.
\end{proof}

\begin{lemma}[Rank is bounded by edge count]
  \label{lem:rigidityMap-finrank-range-le}
  \lean{SimpleGraph.rigidityMap_finrank_range_le}
  \leanok
  \uses{def:rigidityMap}
  If $V$ is finite, then for every framework $p$,
  \[
    \dim_{\R}\,\bigl(\mathrm{range}\,(\mathrm{RigidityMap}\,G\,p)\bigr)
    \;\le\;
    |E(G)|.
  \]
\end{lemma}
\begin{proof}
  \leanok
  The range is a submodule of $E(G) \to \R$, whose dimension is
  $|E(G)|$ by $\mathrm{Module.finrank\_pi}$.
\end{proof}

\subsection{Graph monotonicity}

Generic rigidity is a monotone property in the graph: adding edges
preserves rigidity at any fixed placement.

\begin{lemma}[Infinitesimal rigidity is monotone]
  \label{lem:isInfinitesimallyRigid-mono}
  \lean{SimpleGraph.IsInfinitesimallyRigid.mono}
  \leanok
  \uses{def:isInfinitesimallyRigid, lem:rigidityMap-ker-mono}
  If $G \le G'$ and $p$ is infinitesimally rigid for $G$, then $p$ is
  infinitesimally rigid for $G'$.
\end{lemma}
\begin{proof}
  \leanok
  By \cref{lem:rigidityMap-ker-mono} the kernel of
  $\mathrm{RigidityMap}\,G'\,p$ is contained in that of
  $\mathrm{RigidityMap}\,G\,p$, so its dimension is at most the
  latter's, which is at most $d(d+1)/2$ by hypothesis.
\end{proof}

\begin{lemma}[Generic rigidity is monotone]
  \label{lem:isGenericallyRigid-mono}
  \lean{SimpleGraph.IsGenericallyRigid.mono}
  \leanok
  \uses{def:isGenericallyRigid, lem:isInfinitesimallyRigid-mono}
  If $G \le G'$ and $G$ is generically rigid in dimension $d$, then so
  is $G'$.
\end{lemma}
\begin{proof}
  \leanok
  Apply \cref{lem:isInfinitesimallyRigid-mono} at the same placement
  witness.
\end{proof}

\subsection{The edge-count bound}

\begin{theorem}[Edge-count bound for generically rigid graphs]
  \label{thm:isGenericallyRigid-card-mul-le}
  \lean{SimpleGraph.IsGenericallyRigid.card_mul_le}
  \leanok
  \uses{def:isGenericallyRigid, lem:framework-finrank,
    lem:rigidityMap-finrank-range-le}
  Let $V$ be finite and $G$ generically rigid in dimension $d$. Then
  \[
    d \cdot |V| \;\le\; |E(G)| + \frac{d(d+1)}{2}.
  \]
  Phrased additively to avoid $\N$-subtraction. Specializing at
  $d = 2$ recovers $2 |V| \le |E(G)| + 3$, the necessary half of
  Laman's theorem.
\end{theorem}
\begin{proof}
  \leanok
  Pick an infinitesimally rigid framework $p$ witnessing generic
  rigidity. Rank--nullity for $\mathrm{RigidityMap}\,G\,p$ reads
  \[
    \dim\,\mathrm{Framework}\,V\,d
    \;=\;
    \dim\,(\mathrm{range}) + \dim\,(\ker).
  \]
  By \cref{lem:framework-finrank} the left-hand side equals $d \cdot |V|$;
  by \cref{lem:rigidityMap-finrank-range-le} the rank is at most
  $|E(G)|$; and by definition of infinitesimal rigidity the nullity
  is at most $d(d+1)/2$. Combining yields the claim.
\end{proof}

\subsection{Worked example: \texorpdfstring{$K_2$}{K2} is generically rigid in every dimension}

\begin{theorem}[$K_2$ is generically rigid in every dimension]
  \label{thm:top-fin-two-isGenericallyRigid}
  \lean{SimpleGraph.top_fin_two_isGenericallyRigid}
  \leanok
  \uses{def:isGenericallyRigid, lem:framework-finrank,
    lem:rigidityMap-finrank-range-le, lem:rigidityMap-apply}
  For every $d \in \N$, the complete graph on two vertices is
  generically rigid in dimension $d$.
\end{theorem}
\begin{proof}
  \leanok
  Split on $d$:
  \begin{itemize}
    \item If $d = 0$, the framework space itself is $0$-dimensional,
      so every placement has rigidity-map kernel of dimension $0$
      and the bound $0 \le 0$ holds.
    \item If $d \ge 1$, take $p_0 = 0$ and $p_1 = e_0$ (the first
      standard basis vector). The single-edge image
      $(\mathrm{RigidityMap}\,K_2\,p)(p')(\{0, 1\})
        = \langle p_0 - p_1,\, p'_0 - p'_1 \rangle$
      evaluates at $p' = (e_0, 0)$ to $-1 \ne 0$
      (\cref{lem:rigidityMap-apply}), so the range has dimension at
      least $1$. Rank--nullity (\cref{lem:framework-finrank}) gives
      $2d - 1 \ge \dim\,(\ker)$, and the elementary inequality
      $2d + 1 \le (d+1)(d+2)/2$ for $d \ge 0$ closes the
      kernel-dimension bound to $d(d+1)/2$.
  \end{itemize}
\end{proof}
